\chapter{Chapter05 — AC0 Threshold Parity Separation}
%%% generated by the blueprint pipeline from FABL.Chapter05.AC0ThresholdParitySeparation
\section{Overview}

Declarations of the Lean module \texttt{FABL.Chapter05.AC0ThresholdParitySeparation}.

\section{Declarations}

\begin{theorem}[A quadratic bound by a power of two]
\lean{FABL.two_mul_sq_le_two_pow_of_seven_le}
\label{FABL.two_mul_sq_le_two_pow_of_seven_le}
\leanok
\difficulty{4}
For every integer $m\ge 7$, one has $2m^2\le 2^m$.
\end{theorem}

\begin{theorem}[Active coordinate count bound]
\lean{FABL.activeCoordinateCount_le_dimension}
\label{FABL.activeCoordinateCount_le_dimension}
\leanok
\uses{FABL.two_mul_sq_le_two_pow_of_seven_le}
\difficulty{1}
For every natural number $m\ge 7$, $2m^2\le 2^m$.
\end{theorem}

\begin{definition}[Block coordinate values]
\lean{FABL.blockCoordinateValue}
\label{FABL.blockCoordinateValue}
\leanok
For a natural number $m$, an index $b$ in $\{0,\ldots,m-1\}$, and an index $j$ in
$\{0,\ldots,2m-1\}$, define a natural-number coordinate value as follows: if $j<m$, it
is $bm+j$; otherwise, it is $m^2+bm+(j-m)$.
\end{definition}

\begin{theorem}[Bound on block coordinate values]
\lean{FABL.blockCoordinateValue_lt_active}
\label{FABL.blockCoordinateValue_lt_active}
\leanok
\uses{FABL.blockCoordinateValue}
\difficulty{4}
For every natural number $m$, every index $b<m$, and every index $j<2m$, the block
coordinate value determined by $m,b,j$ is strictly less than $2m^2$.
\end{theorem}

\begin{definition}[Block coordinate index]
\lean{FABL.blockCoordinate}
\label{FABL.blockCoordinate}
\leanok
\uses{FABL.activeCoordinateCount_le_dimension, FABL.blockCoordinateValue, FABL.blockCoordinateValue_lt_active}
For each integer $m \ge 7$, block index $b \in \{0,\ldots,m-1\}$, and position $j \in
\{0,\ldots,2m-1\}$, let $c(m,b,j)$ be the block-coordinate value associated with these
indices. The block coordinate is the index $c(m,b,j)$ in $\{0,\ldots,2^m-1\}$.
\end{definition}

\begin{theorem}[Injectivity of block coordinates]
\lean{FABL.blockCoordinate_injective}
\label{FABL.blockCoordinate_injective}
\leanok
\uses{FABL.blockCoordinate, FABL.blockCoordinateValue}
\difficulty{5}
Let $m\geq 7$, and fix an index $b<m$. The map from indices $j<2m$ to coordinates among
the $2^m$ coordinates, given by the block coordinate construction, is injective.
\end{theorem}

\begin{definition}[Block input]
\lean{FABL.blockInput}
\label{FABL.blockInput}
\leanok
\uses{FABL.blockCoordinate}
For $m\ge 7$, a sign-valued input $x$ with $2^m$ coordinates, and a block index $b<m$,
the block input is the sign-valued vector of length $2m$ whose coordinate at position
$j$ is the coordinate of $x$ selected by the block-coordinate construction for
$(m,b,j)$.
\end{definition}

\begin{definition}[Block minterm]
\lean{FABL.blockMinterm}
\statementsource{Analysis-of-Boolean-Functions-by-Ryan-ODonnell}{pdf-ffbe9c7b3306-p093-b004}
\label{FABL.blockMinterm}
\leanok
\uses{FABL.DNFTerm, FABL.ParityBlockAssignment, FABL.blockCoordinate, FABL.blockCoordinate_injective}
Let $m\ge 7$, choose a block $b\in\{0,\ldots,m-1\}$, and assign a sign in $\{−1,1\}$ to
each of the $2m$ positions in that block's assignment. The associated minterm is the
conjunction of the $2m$ literals formed from those signs and the corresponding block
coordinates among the $2^m$ variables; these coordinates are pairwise distinct.
\end{definition}

\begin{theorem}[Width of a block minterm]
\lean{FABL.width_blockMinterm}
\label{FABL.width_blockMinterm}
\leanok
\uses{FABL.DNFTerm.width, FABL.ParityBlockAssignment, FABL.blockMinterm}
\difficulty{1}
Let $m\ge 7$, choose a block $b\in\{0,\ldots,m-1\}$, and assign a sign to each of its
$2m$ coordinates. The corresponding block minterm has width exactly $2m$.
\end{theorem}

\begin{theorem}[Evaluation of a block minterm]
\lean{FABL.blockMinterm_eval_eq_neg_one_iff}
\label{FABL.blockMinterm_eval_eq_neg_one_iff}
\leanok
\uses{FABL.DNFTerm.eval, FABL.DNFTerm.eval_eq_neg_one_iff, FABL.Literal, FABL.ParityBlockAssignment, FABL.blockCoordinate, FABL.blockInput, FABL.blockMinterm}
\difficulty{4}
Let $m\geq 7$, let $b<m$, let $a$ assign a sign to each of the $2m$ positions, and let
$x$ be a sign vector with $2^m$ coordinates. The block minterm determined by $m,b,a$
evaluates to $-1$ on $x$ if and only if $a$ agrees at every position with the
restriction of $x$ to the block coordinates indexed by $b$.
\end{theorem}

\begin{definition}[Block layer size]
\lean{FABL.blockLayerSize}
\label{FABL.blockLayerSize}
\leanok
For a natural number $m$, the block layer size is $m\,2^{2m}$.
\end{definition}

\begin{definition}[Block assignment index equivalence]
\lean{FABL.blockAssignmentIndexEquiv}
\label{FABL.blockAssignmentIndexEquiv}
\leanok
\uses{FABL.ParityBlockAssignment, FABL.blockLayerSize, FABL.parityBlockAssignmentEquiv}
For each natural number $m$, there is a bijection between the indices
$0,\ldots,m2^{2m}-1$ and pairs consisting of a block index below $m$ and an assignment
of a sign to each of the $2m$ positions.
\end{definition}

\begin{definition}[Block first layer]
\lean{FABL.blockFirstLayer}
\label{FABL.blockFirstLayer}
\leanok
\uses{FABL.DNFTerm, FABL.blockAssignmentIndexEquiv, FABL.blockLayerSize, FABL.blockMinterm}
For $m\ge 7$, the first layer is the list of DNF terms obtained by ranging over all
pairs of a block index below $m$ and a sign assignment on $2m$ positions, and taking the
block minterm associated with each pair. The list has $m2^{2m}$ entries.
\end{definition}

\begin{theorem}[Length of the first block layer]
\lean{FABL.length_blockFirstLayer}
\label{FABL.length_blockFirstLayer}
\leanok
\uses{FABL.blockFirstLayer, FABL.blockLayerSize}
\difficulty{1}
For every natural number $m$ with $7\le m$, the first block layer contains exactly
$m\,2^{2m}$ terms.
\end{theorem}

\begin{theorem}[First-layer gate characterization]
\lean{FABL.eval_blockFirstLayer}
\label{FABL.eval_blockFirstLayer}
\leanok
\uses{FABL.CircuitGate.evalTerm, FABL.blockAssignmentIndexEquiv, FABL.blockFirstLayer, FABL.blockInput, FABL.blockLayerSize, FABL.blockMinterm, FABL.blockMinterm_eval_eq_neg_one_iff, FABL.length_blockFirstLayer}
\difficulty{3}
Let $m\ge 7$, let $x$ be a sign string with $2^m$ coordinates, and let $i$ index one of
the $m2^{2m}$ first-layer gates. Decompose $i$ through the specified indexing bijection
into a block index $b\in\{0,\ldots,m-1\}$ and a sign assignment $a$ on $2m$ positions.
The corresponding AND gate in the first layer outputs $-1$ exactly when $a$ agrees at
every position with the restriction of $x$ to the coordinates selected for block $b$.
\end{theorem}

\begin{definition}[Block inner product value]
\lean{FABL.blockInnerProductValue}
\label{FABL.blockInnerProductValue}
\leanok
\uses{FABL.Sign, FABL.blockInput, FABL.innerProductModTwoBoolean}
For $m\ge 7$, a sign-valued input $x$ with $2^m$ coordinates, and a block index $b<m$,
the block inner product value is the parity inner product modulo two of the length-$2m$
block input selected by $b$.
\end{definition}

\begin{definition}[Block output pattern]
\lean{FABL.blockOutputPattern}
\label{FABL.blockOutputPattern}
\leanok
\uses{FABL.ParityBlockAssignment, FABL.blockInnerProductValue}
For $m\ge 7$ and a sign-valued input $x$ with $2^m$ coordinates, the block output
pattern assigns to each block index $b<m$ the block inner product value of $x$ on that
block.
\end{definition}

\begin{definition}[Threshold parity target]
\lean{FABL.ac0ThresholdParityTarget}
\statementsource{mit-18405-razborov-smolensky}{
  pdf-d3f464808105-p002-b001,
  pdf-d3f464808105-p007-b002
}
\label{FABL.ac0ThresholdParityTarget}
\leanok
\uses{FABL.BooleanFunction, FABL.activeCoordinateCount_le_dimension, FABL.blockCoordinateValue, FABL.blockCoordinateValue_lt_active, FABL.innerProductModTwoBoolean}
Let $m\ge 7$ and consider inputs with $2^m$ signs. For each $b\in\{0,\ldots,m-1\}$, take
the parity, modulo two, of the $m$ input signs selected by the block-coordinate rule for
$b$; the target function is the product of these $m$ parity values.
\end{definition}

\begin{theorem}[Parity target as a product of block outputs]
\lean{FABL.ac0ThresholdParityTarget_apply}
\label{FABL.ac0ThresholdParityTarget_apply}
\leanok
\uses{FABL.ac0ThresholdParityTarget, FABL.blockOutputPattern, FABL.parityBlockProduct}
\difficulty{3}
Let $m\ge 7$, and let $x\in\{-1,1\}^{2^m}$. For each $b\in\{0,\ldots,m-1\}$, take the
$2m$ coordinates of $x$ indexed by the two blocks specified by $b$, and let
$q_b(x)\in\{-1,1\}$ be the sign-valued mod-two inner product of those blocks. Then the
parity target evaluated at $x$ is $\prod_{b=0}^{m-1}q_b(x)$.
\end{theorem}

\begin{definition}[Clause node inputs]
\lean{FABL.clauseNodeInputs}
\label{FABL.clauseNodeInputs}
\leanok
\uses{FABL.DNFTerm, FABL.blockAssignmentIndexEquiv, FABL.blockLayerSize, FABL.innerProductModTwoBoolean}
For a natural number $m$ and a DNF term $C$ on $m$ coordinates, the clause node inputs
are those indices in a set of size $m2^{2m}$ whose associated block index and sign
assignment match the index and required sign of at least one literal of $C$; the match
requires equality of the block index with the literal’s coordinate and equality of the
parity inner product of the assignment with the literal’s required sign.
\end{definition}

\begin{definition}[Block middle layer]
\lean{FABL.blockMiddleLayer}
\label{FABL.blockMiddleLayer}
\leanok
\uses{FABL.blockLayerSize, FABL.clauseNodeInputs, FABL.parityCNF}
For each natural number $m$, form the list of input-coordinate sets associated to the
clauses of the parity CNF on $m$ variables. Its entries are finite subsets of a
coordinate set of size $m2^{2m}$.
\end{definition}

\begin{theorem}[Length of the middle block layer]
\lean{FABL.length_blockMiddleLayer}
\label{FABL.length_blockMiddleLayer}
\leanok
\uses{FABL.blockMiddleLayer, FABL.size_parityCNF}
\difficulty{3}
For every natural number $m$ with $7\le m$, the middle block layer contains exactly
$2^{m-1}$ clauses.
\end{theorem}

\begin{definition}[Threshold parity circuit]
\lean{FABL.ac0ThresholdParityCircuit}
\statementsource{Analysis-of-Boolean-Functions-by-Ryan-ODonnell}{
  pdf-ffbe9c7b3306-p103-b002,
  pdf-ffbe9c7b3306-p104-b001,
  pdf-ffbe9c7b3306-p116-b001,
  pdf-ffbe9c7b3306-p116-b002
}
\label{FABL.ac0ThresholdParityCircuit}
\leanok
\uses{FABL.CircuitTail, FABL.DepthCircuit, FABL.blockFirstLayer, FABL.blockMiddleLayer, FABL.length_blockFirstLayer}
Let $m\ge 7$. The threshold parity circuit on $2^m$ inputs has a first layer of AND
gates given by the block minterms, followed by the specified middle layer and a
singleton output layer.
\end{definition}

\begin{theorem}[Depth of the threshold parity circuit]
\lean{FABL.depth_ac0ThresholdParityCircuit}
\label{FABL.depth_ac0ThresholdParityCircuit}
\leanok
\uses{FABL.CircuitTail, FABL.CircuitTail.layerCount, FABL.DepthCircuit.depth, FABL.ac0ThresholdParityCircuit, FABL.blockMiddleLayer, FABL.circuitTail_layerCount_cast, FABL.length_blockFirstLayer}
\difficulty{3}
For every integer $m\ge 7$, the constructed circuit on $2^m$ inputs has depth exactly
$3$.
\end{theorem}

\begin{theorem}[Size bound for a parity circuit]
\lean{FABL.size_ac0ThresholdParityCircuit}
\label{FABL.size_ac0ThresholdParityCircuit}
\leanok
\uses{FABL.CircuitTail, FABL.CircuitTail.internalNodeCount, FABL.DepthCircuit.size, FABL.ac0ThresholdParityCircuit, FABL.blockFirstLayer, FABL.blockLayerSize, FABL.blockMiddleLayer, FABL.circuitTail_internalNodeCount_cast, FABL.length_blockFirstLayer, FABL.length_blockMiddleLayer}
\difficulty{4}
For every integer $m\\ge 7$, the circuit constructed from blocks of $2m$ parity
coordinates and the parity conjunctive normal form on $m$ variables has size at most
$m2^{2m}+2^m$.
\end{theorem}

\begin{theorem}[Width of the first DNF layer]
\lean{FABL.blockFirstLayer_width}
\label{FABL.blockFirstLayer_width}
\leanok
\uses{FABL.DNFFormula.width, FABL.DNFFormula.width_le_iff, FABL.DNFFormula.width_le_of_mem, FABL.blockAssignmentIndexEquiv, FABL.blockFirstLayer, FABL.blockLayerSize, FABL.blockMinterm, FABL.length_blockFirstLayer, FABL.width_blockMinterm}
\difficulty{4}
For every natural number $m\geq 7$, the width of the DNF formula whose terms are the
first block layer is $2m$.
\end{theorem}

\begin{theorem}[Width of the Threshold-Parity Circuit]
\lean{FABL.width_ac0ThresholdParityCircuit}
\label{FABL.width_ac0ThresholdParityCircuit}
\leanok
\uses{FABL.DepthCircuit.width, FABL.ac0ThresholdParityCircuit, FABL.blockFirstLayer_width}
\difficulty{3}
For every integer $m\geq 7$, consider the circuit on $2^m$ sign-valued inputs whose
first layer consists of one minterm for each pair comprising an index
$b\in\{0,\ldots,m-1\}$ and a sign assignment to $2m$ positions, with those positions
mapped to the corresponding block coordinates. The maximum number of literals in a
first-layer term, which is the circuit's width, is exactly $2m$.
\end{theorem}

\begin{theorem}[Evaluation of clause node inputs]
\lean{FABL.eval_clauseNodeInputs}
\label{FABL.eval_clauseNodeInputs}
\leanok
\uses{FABL.CNFFormula.clauseEval, FABL.CNFFormula.sign_eq_of_ne_neg, FABL.CircuitGate.evalFinset, FABL.CircuitGate.evalFinset_or_eq_neg_one_iff, FABL.CircuitGate.evalTerm, FABL.DNFTerm, FABL.Literal.eval, FABL.blockAssignmentIndexEquiv, FABL.blockFirstLayer, FABL.blockInnerProductValue, FABL.blockInput, FABL.blockLayerSize, FABL.blockOutputPattern, FABL.clauseNodeInputs, FABL.eval_blockFirstLayer, FABL.length_blockFirstLayer}
\difficulty{5}
Let $m	ext{ and }x$ be such that $7le m$ and $x$ is a sign assignment on $2^m$
coordinates, and let $C$ be a conjunction of literals on $m$ coordinates with distinct
indices. For each first-layer node, evaluate its conjunction of literals on $x$; then
take the disjunction over the nodes whose block index is the index of some literal of
$C$ and whose parity output equals that literal's required sign. The resulting sign
equals the clause evaluation of $C$ on the assignment whose value at each block is the
parity output of that block.
\end{theorem}

\begin{theorem}[Evaluation of an individual middle gate]
\lean{FABL.eval_blockMiddleLayer_get}
\label{FABL.eval_blockMiddleLayer_get}
\leanok
\uses{FABL.CNFFormula.clauseEval, FABL.CircuitGate.evalFinset, FABL.CircuitGate.evalTerm, FABL.blockFirstLayer, FABL.blockLayerSize, FABL.blockMiddleLayer, FABL.blockOutputPattern, FABL.clauseNodeInputs, FABL.eval_clauseNodeInputs, FABL.length_blockFirstLayer, FABL.parityCNF}
\difficulty{1}
Let $m\ge 7$, let $x$ be a sign string with $2^m$ coordinates, and let $i$ index a
clause in the middle layer. The OR gate on the inputs specified for clause $i$, fed by
the corresponding first-layer AND outputs, has the same value as that clause evaluated
on the length-$m$ sign pattern whose entry at block $b$ is the parity sign of the
two-block restriction of $x$ selected for $b$.
\end{theorem}

\begin{theorem}[Evaluation of the middle circuit layer]
\lean{FABL.eval_blockMiddleLayer}
\label{FABL.eval_blockMiddleLayer}
\leanok
\uses{FABL.CNFFormula.clauseEval, FABL.CNFFormula.eval, FABL.CNFFormula.sign_eq_of_ne_neg, FABL.CircuitGate.evalFinset, FABL.CircuitGate.evalFinset_and_eq_neg_one_iff, FABL.CircuitGate.evalTerm, FABL.blockFirstLayer, FABL.blockLayerSize, FABL.blockMiddleLayer, FABL.blockOutputPattern, FABL.eval_blockMiddleLayer_get, FABL.length_blockFirstLayer, FABL.parityCNF}
\difficulty{5}
Let $m\ge 7$ and let $x$ be a sign string with $2^m$ coordinates. Evaluate the middle
layer by feeding each of its OR gates the outputs of the corresponding first-layer AND
gates, then take the AND of all those OR outputs. The resulting sign equals the value of
the parity conjunctive normal form on the length-$m$ sign pattern whose entry at block
$b$ is the parity sign of the two-block restriction of $x$ selected for $b$.
\end{theorem}

\begin{theorem}[Evaluation of the parity circuit]
\lean{FABL.eval_ac0ThresholdParityCircuit}
\label{FABL.eval_ac0ThresholdParityCircuit}
\leanok
\uses{FABL.CNFFormula.toBooleanFunction, FABL.CircuitGate.dual_or, FABL.CircuitGate.evalTerm, FABL.CircuitTail, FABL.CircuitTail.eval, FABL.DepthCircuit.eval, FABL.ac0ThresholdParityCircuit, FABL.ac0ThresholdParityTarget, FABL.ac0ThresholdParityTarget_apply, FABL.blockFirstLayer, FABL.blockMiddleLayer, FABL.blockOutputPattern, FABL.circuitTail_eval_cast, FABL.eval_blockMiddleLayer, FABL.length_blockFirstLayer, FABL.parityBlockProduct, FABL.parityCNF, FABL.parityCNF_toBooleanFunction, FABL.parityFunction}
\difficulty{3}
Let $m\ge 7$ and let $x$ be a sign string with $2^m$ coordinates. For each
$b\in\{0,\ldots,m-1\}$, form a string of length $2m$ by taking from $x$ the coordinates
indexed, for $0\le j<2m$, by $bm+j$ when $j<m$, and by $m^2+bm+(j-m)$ when $j\ge m$. The
constructed circuit with AND and OR gates evaluates on $x$ to the sign obtained by
multiplying, over all such $b$, the parity signs of these two-block strings.
\end{theorem}

\begin{theorem}[Polynomial bound for circuit size]
\lean{FABL.ac0ThresholdParityCircuit_size_le_cube}
\label{FABL.ac0ThresholdParityCircuit_size_le_cube}
\leanok
\difficulty{4}
For every natural number $m$, $m\,2^{2m}+2^m\le (2^m)^3$.
\end{theorem}

\begin{theorem}[Circuit for a parity product]
\lean{FABL.hasDepthCircuit_ac0ThresholdParityTarget}
\label{FABL.hasDepthCircuit_ac0ThresholdParityTarget}
\leanok
\uses{FABL.DepthCircuit.HasDepthCircuit, FABL.ac0ThresholdParityCircuit, FABL.ac0ThresholdParityCircuit_size_le_cube, FABL.ac0ThresholdParityTarget, FABL.depth_ac0ThresholdParityCircuit, FABL.eval_ac0ThresholdParityCircuit, FABL.size_ac0ThresholdParityCircuit, FABL.width_ac0ThresholdParityCircuit}
\difficulty{3}
Let $m$ be a natural number with $7\le m$. On the sign cube with $2^m$ coordinates, form
$m$ pairs of blocks, each block having $m$ coordinates: for each $b=0,\ldots,m-1$, pair
the coordinates $bm,\ldots,bm+m-1$ with the coordinates $m^2+bm,\ldots,m^2+bm+m-1$.
Define the target function as the product, over these $m$ pairs, of the sign encoding of
the dot product modulo $2$ of the two blocks. This function is computed by a depth-three
circuit with at most $(2^m)^3$ internal gates and width at most $2m$.
\end{theorem}

\begin{definition}[Reindexing a sign cube]
\lean{FABL.reindexSignCube}
\label{FABL.reindexSignCube}
\leanok
\uses{FABL.Sign}
Let $\kappa$ and $\iota$ be types, let $e:\kappa\simeq\iota$ be a bijection, and let the
sign set be $\{1,-1\}$. Reindexing functions by $e$ gives a bijection from sign-valued
functions on $\kappa$ to sign-valued functions on $\iota$, sending $x$ to the function
$i\mapsto x(e^{-1}(i))$.
\end{definition}

\begin{theorem}[Fourier coefficients under coordinate reindexing]
\lean{FABL.fourierCoeff_reindexSignCube}
\label{FABL.fourierCoeff_reindexSignCube}
\leanok
\uses{FABL.Sign, FABL.fourierCoeff, FABL.indexedFourierCoeff, FABL.indexedMonomial, FABL.monomial, FABL.reindexSignCube}
\difficulty{4}
Let $k$ be a nonnegative integer, let $\iota$ be a finite index set, and let
$e:[k]\to\iota$ be a bijection. For any real-valued function $f$ on sign assignments to
$\iota$ and any subset $S\subseteq[k]$, the Fourier coefficient at $S$ of the function
obtained by reindexing inputs through $e$ equals the Fourier coefficient of $f$ at the
image $e(S)$.
\end{theorem}

\begin{definition}[Reindexed restriction]
\lean{FABL.reindexedRestriction}
\statementsource{Analysis-of-Boolean-Functions-by-Ryan-ODonnell}{pdf-ffbe9c7b3306-p075-b002}
\label{FABL.reindexedRestriction}
\leanok
\uses{FABL.FixedSignCube, FABL.reindexSignCube, FABL.signRestriction}
Let $n$ and $k$ be natural numbers, let $p$ be a real-valued function on the sign cube
of dimension $n$, let $J\subseteq\mathrm{Fin}(n)$, and fix signs on the coordinates
outside $J$. Given a bijection from $\mathrm{Fin}(k)$ to $J$, define a function on the
sign cube of dimension $k$ by restricting $p$ to those fixed outside signs and
reindexing the remaining coordinates using the bijection.
\end{definition}

\begin{theorem}[Fourier coefficients of reindexed restrictions]
\lean{FABL.fourierCoeff_reindexedRestriction}
\label{FABL.fourierCoeff_reindexedRestriction}
\leanok
\uses{FABL.FixedSignCube, FABL.fourierCoeff, FABL.fourierCoeff_reindexSignCube, FABL.reindexedRestriction, FABL.restrictionFourierCoeff, FABL.signRestriction}
\difficulty{1}
Let $n,k$ be nonnegative integers, let $J\subseteq[n]$, fix a sign assignment $z$ on the
coordinates outside $J$, and let $e:[k]\to J$ be a bijection. For any real-valued
function $p$ on the sign cube of dimension $n$ and any $S\subseteq[k]$, the Fourier
coefficient at $S$ of the restriction with free coordinates reindexed by $e$ equals the
Fourier coefficient, on the restricted cube, at $e(S)$ with the fixed assignment $z$.
\end{theorem}

\begin{theorem}[Ambient frequency from restricted Fourier support]
\lean{FABL.exists_ambientFrequency_of_reindexedRestriction_memSupport}
\proofsource{Analysis-of-Boolean-Functions-by-Ryan-ODonnell}{
  pdf-ffbe9c7b3306-p076-b003,
  pdf-ffbe9c7b3306-p076-b004
}
\label{FABL.exists_ambientFrequency_of_reindexedRestriction_memSupport}
\leanok
\uses{FABL.FixedIndex, FABL.FixedSignCube, FABL.fourierCoeff, FABL.fourierCoeff_reindexedRestriction, FABL.fourierSupport, FABL.liftFixedFrequency, FABL.liftFreeFrequency, FABL.mem_fourierSupport, FABL.reindexedRestriction, FABL.restrictionFourierCoeff_eq_sum}
\difficulty{4}
Let $n,k$ be nonnegative integers, let $J\subseteq\{0,\ldots,n-1\}$, let $p:\{
-1,1\}^n\to\mathbb R$, and fix signs $z$ on the coordinates outside $J$. Let
$e:\{0,\ldots,k-1\}\to J$ be a bijection, and let $S\subseteq\{0,\ldots,k-1\}$. If $S$
belongs to the Fourier support of the restriction of $p$ to the free coordinates,
reindexed by $e$, then there is a set $T$ of coordinates outside $J$ such that the
Fourier coefficient of $p$ at the union of the ambient lift of $e(S)$ and the lift of
$T$ is nonzero.
\end{theorem}

\begin{theorem}[Free part of a split frequency set]
\lean{FABL.freeFrequencyPart_lift_union}
\label{FABL.freeFrequencyPart_lift_union}
\leanok
\uses{FABL.FixedIndex, FABL.freeFrequencyPart, FABL.liftFixedFrequency, FABL.liftFreeFrequency}
\difficulty{3}
Let $n$ be a natural number and let $J$ be a subset of the coordinates
$\{0,\ldots,n-1\}$. For any subset $S$ of $J$ and any subset $T$ of its complement, the
part inside $J$ of the union of $S$ and $T$ is $S$.
\end{theorem}

\begin{theorem}[Sparsity under restriction and reindexing]
\lean{FABL.polynomialSparsity_reindexedRestriction_le}
\label{FABL.polynomialSparsity_reindexedRestriction_le}
\leanok
\uses{FABL.FixedIndex, FABL.FixedSignCube, FABL.exists_ambientFrequency_of_reindexedRestriction_memSupport, FABL.fourierCoeff, FABL.fourierSupport, FABL.freeFrequencyPart, FABL.freeFrequencyPart_lift_union, FABL.liftFixedFrequency, FABL.liftFreeFrequency, FABL.mem_fourierSupport, FABL.polynomialSparsity, FABL.reindexedRestriction}
\difficulty{4}
Let $n,k$ be natural numbers, let $p$ be a real-valued function on the $n$-dimensional
sign cube, let $J$ be a set of coordinates, let $z$ assign a sign to each coordinate
outside $J$, and let $e$ be a bijection from the $k$ coordinates to $J$. The number of
nonzero Fourier coefficients of the restriction of $p$ obtained by fixing the
coordinates outside $J$ to $z$ and reindexing the remaining coordinates by $e$ is at
most the number of nonzero Fourier coefficients of $p$.
\end{theorem}

\begin{definition}[Active coordinates]
\lean{FABL.ac0ActiveCoordinates}
\label{FABL.ac0ActiveCoordinates}
\leanok
For each natural number $m$, the active coordinates are those indices $i$ in the set of
integers from $0$ through $2^m-1$ satisfying $i<m^2+m^2$.
\end{definition}

\begin{definition}[Active coordinate bijection]
\lean{FABL.ac0ActiveCoordinateEquiv}
\label{FABL.ac0ActiveCoordinateEquiv}
\leanok
\uses{FABL.ac0ActiveCoordinates, FABL.activeCoordinateCount_le_dimension}
Let $m\ge 7$. The active coordinates among the indices from $0$ through $2^m-1$ are in
bijection with the indices from $0$ through $m^2+m^2-1$, by identifying each index with
the same integer.
\end{definition}

\begin{definition}[Inactive coordinate assignment]
\lean{FABL.ac0InactiveAssignment}
\label{FABL.ac0InactiveAssignment}
\leanok
\uses{FABL.FixedSignCube, FABL.ac0ActiveCoordinates}
For each natural number $m$, assign the sign $+1$ to every coordinate outside the
active-coordinate set.
\end{definition}

\begin{theorem}[Recovery of active sign coordinates]
\lean{FABL.combine_ac0ActiveCoordinates}
\label{FABL.combine_ac0ActiveCoordinates}
\leanok
\uses{FABL.ac0ActiveCoordinateEquiv, FABL.ac0ActiveCoordinates, FABL.ac0InactiveAssignment, FABL.activeCoordinateCount_le_dimension, FABL.combineSignCube, FABL.combineSignCube_apply_free, FABL.reindexSignCube}
\difficulty{3}
Let $m\ge 7$, let $y$ be a sign vector of length $2m^2$, and let $i$ be an index in that
vector. In a sign cube of dimension $2^m$, take the active coordinates to be those with
index less than $2m^2$, place the entries of $y$ on those coordinates in their natural
order, and assign $+1$ to every other coordinate. The resulting vector has value $y(i)$
at the coordinate indexed by $i$.
\end{theorem}

\begin{theorem}[Sign encoding of a finite sum]
\lean{FABL.signEncode_sum}
\label{FABL.signEncode_sum}
\leanok
\uses{FABL.signEncode, FABL.signEncode_add, FABL.𝔽₂}
\difficulty{4}
Let $\iota$ be a type, let $S$ be a finite subset of $\iota$, and let
$a_i\in\mathbb{F}_2$ for each $i\in\iota$. Encode each $b\in\mathbb{F}_2$ by the unit
$(-1)^b$ in the integers. Then the encoding of the sum of the $a_i$ over $S$ equals the
product of their encodings: $(-1)^{\sum_{i\in S}a_i}=\prod_{i\in S}(-1)^{a_i}$, with the
sum taken in $\mathbb{F}_2$.
\end{theorem}

\begin{theorem}[Block product of mod-two inner products]
\lean{FABL.prod_innerProductModTwoBoolean_blocks}
\label{FABL.prod_innerProductModTwoBoolean_blocks}
\leanok
\uses{FABL.binaryCubeSignEquiv, FABL.f₂DotProduct, FABL.innerProductModTwoBoolean, FABL.innerProductModTwoBoolean_binaryCubeSignEquiv_joinF₂CubeBlocks, FABL.joinF₂CubeBlocks, FABL.signEncode_sum}
\difficulty{4}
Let $m$ be a natural number and let $a,c$ be arrays over the two-element field indexed
by pairs of elements of an $m$-element set. For each row $b$, take the sign
$(-1)^{	ext{row}_a(b)oldsymbol{3}	ext{row}_c(b)}$, where the dot product is computed
modulo $2$. The product of these $m$ signs equals $(-1)^{aoldsymbol{3}c}$, with the
dot product on the right taken over all $m^2$ coordinates modulo $2$.
\end{theorem}

\begin{theorem}[Active restriction of the parity target]
\lean{FABL.ac0ThresholdParityTarget_activeRestriction}
\label{FABL.ac0ThresholdParityTarget_activeRestriction}
\leanok
\uses{FABL.ac0ActiveCoordinateEquiv, FABL.ac0ActiveCoordinates, FABL.ac0InactiveAssignment, FABL.ac0ThresholdParityTarget, FABL.activeCoordinateCount_le_dimension, FABL.binaryCubeSignEquiv, FABL.blockCoordinateValue, FABL.blockCoordinateValue_lt_active, FABL.combineSignCube, FABL.combine_ac0ActiveCoordinates, FABL.f₂CubeBlockEquiv, FABL.innerProductModTwoBoolean, FABL.joinF₂CubeBlocks, FABL.prod_innerProductModTwoBoolean_blocks, FABL.reindexSignCube}
\difficulty{6}
Let $m\ge 7$, and let $y$ be a sign vector of length $2m^2$. Form an input of length
$2^m$ by placing the coordinates of $y$ on the active coordinates, in the specified
reindexing, and assigning $+1$ to every inactive coordinate. The target function on this
input equals the sign-valued mod-two inner-product function on the two length-$m^2$
blocks of $y$.
\end{theorem}

\begin{theorem}[Sparsity of the constructed threshold parity target]
\lean{FABL.pow_two_sq_le_polynomialSparsity_ac0ThresholdParityTarget}
\label{FABL.pow_two_sq_le_polynomialSparsity_ac0ThresholdParityTarget}
\leanok
\uses{FABL.IsPolynomialThresholdRepresentation, FABL.ac0ActiveCoordinateEquiv, FABL.ac0ActiveCoordinates, FABL.ac0InactiveAssignment, FABL.ac0ThresholdParityTarget, FABL.ac0ThresholdParityTarget_activeRestriction, FABL.combineSignCube, FABL.innerProductModTwoBoolean, FABL.polynomialSparsity, FABL.polynomialSparsity_reindexedRestriction_le, FABL.pow_two_le_polynomialSparsity_innerProductModTwo_of_pos, FABL.reindexSignCube, FABL.reindexedRestriction, FABL.signRestriction}
\difficulty{4}
Let $m\ge 7$, and define a Boolean function on $2^m$ signs by taking the product over
$b=0,\ldots,m-1$ of the signs encoding the mod-two inner products of the two length-$m$
blocks selected from coordinates indexed by $b$ in the first and second halves of the
coordinate range. If a real-valued function on this sign cube represents the target by
threshold signs, then it has at least $2^{m^2}$ nonzero Fourier coefficients.
\end{theorem}

\begin{theorem}[Threshold parity size for the constructed target]
\lean{FABL.pow_two_pow_log_le_thresholdOfParitiesSize_ac0ThresholdParityTarget}
\label{FABL.pow_two_pow_log_le_thresholdOfParitiesSize_ac0ThresholdParityTarget}
\leanok
\uses{FABL.IsThresholdOfParities, FABL.ac0ThresholdParityTarget, FABL.polynomialSparsity, FABL.pow_two_sq_le_polynomialSparsity_ac0ThresholdParityTarget}
\difficulty{4}
Let $m\ge 7$ and let $s$ be a nonnegative integer. Define a Boolean function on $2^m$
signs by taking the product, over $b=0,\ldots,m-1$, of the signs encoding the mod-two
inner products of the two length-$m$ blocks selected from the coordinates indexed by $b$
in the first and second halves of the coordinate range. If this function has a
polynomial threshold representation with at most $s$ nonzero Fourier coefficients, then
$(2^m)^{\lfloor\log_2(2^m)\rfloor}\le s$.
\end{theorem}
